\chapter{Локальный предел часового квантового марковского процесса}
\label{ch:v8-local-observable-clocked-qms-limit-and-time-anchor-gate}

Предыдущая глава отделила глобальную норму бесконечного конвейера от
локальных наблюдаемых. Теперь совместим две аппроксимации, которые до сих
пор рассматривались раздельно: слабое столкновительное приближение
непрерывной полугруппы и конечную точность автономных часов.

\section{Сложение двух ошибок}

Пусть $\Phi_{u/n}$ --- один слабый шаг полного 42-скачкового процесса.
В конечномерной системе правило Чернова даёт на фиксированном интервале
$u$ оценку
\begin{equation}
 \left\|\Phi_{u/n}^{,n}-e^{u\mathcal L_{42}}\right\|
 \leq \frac{C_u}{n}.
 \label{eq:v8-clocked-qms-collision-error}
\end{equation}
Если автономные часы реализуют каждый шаг с ошибкой
$\varepsilon_d\leq Ae^{-cd}$, то телескопическая оценка $n$ композиций
даёт
\begin{equation}
 \boxed{
 \left\|\widetilde\Phi_{u/n,d}^{,n}-e^{u\mathcal L_{42}}\right\|
 \leq \frac{C_u}{n}+nAe^{-cd}.}
 \label{eq:v8-clocked-qms-total-error}
\end{equation}
Первый член измеряет дискретизацию столкновений, второй --- накопленную
ошибку часового управления.

\section{Совместный предел}

Для любого $\alpha>0$ выберем
\begin{equation}
 d_n\geq \frac{1+\alpha}{c}\log n.
 \label{eq:v8-clocked-qms-dimension-schedule}
\end{equation}
Тогда
\begin{equation}
 nAe^{-cd_n}\leq \frac{A}{n^\alpha},
 \qquad
 \left\|\widetilde\Phi_{u/n,d_n}^{,n}
 -e^{u\mathcal L_{42}}\right\|
 \leq \frac{C_u}{n}+\frac{A}{n^\alpha}longrightarrow0.
 \label{eq:v8-clocked-qms-joint-limit}
\end{equation}

\begin{theorem}[Локальный автономный предел непрерывного шума]
На наблюдаемых конечной системы и на любом фиксированном цилиндре цепи
вспомогательных систем полный 42-скачковый квантовый марковский процесс
получается совместным пределом слабых столкновений и автономных часов.
Для этого достаточно логарифмического роста размерности часов по числу
столкновений.
\end{theorem}

Утверждение относится к редуцированным наблюдаемым. Оно не означает
сходимости чистых состояний всей бесконечной цепи в глобальной норме и не
создаёт внешнего пошагового сброса: свежие вспомогательные системы уже
поставляет предварительно подготовленный конвейер.

\section{Почему часы всё ещё не создают секунду}

Обозначим частоту часов через $\Omega$, а физическую интенсивность
диссипативного процесса через $\Gamma$. Наблюдаемая динамика зависит от
безразмерных комбинаций
\begin{equation}
 \Omega t_{\rm phys},\qquad \Gamma t_{\rm phys}.
\end{equation}
Преобразование
\begin{equation}
 (\Omega,\Gamma,t_{\rm phys})\longmapsto
 (\lambda\Omega,\lambda\Gamma,t_{\rm phys}/\lambda),
 \qquad \lambda>0,
 \label{eq:v8-clocked-qms-common-scale-orbit}
\end{equation}
сохраняет обе комбинации. Часы могут определить отношение
$\Gamma/\Omega$, если взаимодействие типизированно связывает их с системой,
но общий размерный множитель остаётся свободным.

\begin{nogocriterion}[Автономность не равна абсолютной калибровке времени]
Совместный предел \eqref{eq:v8-clocked-qms-joint-limit} выводит непрерывный
безразмерный процесс без внешнего переключателя. Он не выбирает секунду:
для этого необходим хотя бы один независимо фиксированный энергетический
или скоростной масштаб и доказанная связь этого масштаба с $\Omega$ либо
$\Gamma$.
\end{nogocriterion}

\begin{protocol}
\begin{enumerate}
 \item ошибка столкновительного предела: \textbf{$C_u/n$};
 \item накопленная ошибка часов: \textbf{$nAe^{-cd}$};
 \item совместная оценка: \textbf{$C_u/n+nAe^{-cd}$};
 \item достаточный рост часов: \textbf{$d_n\geq(1+\alpha)c^{-1}\log n$};
 \item редуцированный непрерывный предел: \textbf{получен};
 \item внешний пошаговый сброс: \textbf{не требуется};
 \item глобальная норма бесконечной цепи: \textbf{не заявляется};
 \item относительная частота $\Gamma/\Omega$: \textbf{может быть измерена};
 \item абсолютная секунда: \textbf{не выведена};
 \item следующий гейт: \textbf{типизированный мост энергии часов к
 интенсивности полного шумового процесса}.
\end{enumerate}
\end{protocol}

Машинный сертификат сохранён в
\path{s2t_v8_local_observable_clocked_qms_limit_and_time_anchor_gate_results.json}.

\begin{thebibliography}{9}
\bibitem{v8-clocked-qms-attal-pautrat}
S.~Attal, Y.~Pautrat,
\emph{From repeated to continuous quantum interactions},
Ann. Henri Poincar\'e \textbf{7} (2006), 59--104;
arXiv:math-ph/0311002.
\bibitem{v8-clocked-qms-attal-joye}
S.~Attal, A.~Joye,
\emph{Weak coupling and continuous limits for repeated quantum interactions},
J. Stat. Phys. \textbf{126} (2007), 1241--1283;
arXiv:math-ph/0501012.
\bibitem{v8-clocked-qms-woods}
M.~P.~Woods, R.~Silva, J.~Oppenheim,
\emph{Autonomous quantum machines and the finite sized Quasi-Ideal clock},
Ann. Henri Poincar\'e \textbf{20} (2019), 125--218;
arXiv:1607.04591.
\end{thebibliography}